%% GENERATED by generate_papers.py -- edit the bank_*.py files, not this file.
\documentclass[11pt,a4paper]{article}
\usepackage[a4paper, top=20mm, bottom=20mm, left=18mm, right=18mm, headheight=14pt]{geometry}
\usepackage[T1]{fontenc}
\usepackage{lmodern}
\usepackage{amsmath, amssymb, mathtools}
\usepackage{braket}
\usepackage{array, tabularx, booktabs, multicol}
\usepackage{enumitem}
\usepackage[dvipsnames]{xcolor}
\usepackage[most]{tcolorbox}
\usepackage{needspace}
\usepackage{fancyhdr, lastpage}
\usepackage{pifont}
\usepackage{tikz}
\usepackage{graphicx}
\usepackage{eso-pic}
% SILICO FELLER watermark: official logo, centred, ~8 % opacity, behind all content, on every page
\AddToShipoutPictureBG{\AtPageCenter{\makebox(0,0){\tikz\node[opacity=0.08, inner sep=0pt]{\includegraphics[width=0.72\paperwidth]{../../../assets/branding/silico_feller_logo.png}};}}}
\usepackage[hidelinks]{hyperref}
\newcolumntype{L}{>{\raggedright\arraybackslash}X}
\definecolor{brand}{HTML}{1F4E79}
\definecolor{okgreen}{HTML}{1B7F3B}
\setlength{\parindent}{0pt}
\setlength{\parskip}{3pt}
\newcommand{\optlabel}[1]{\textbf{(#1)}}
\newcommand{\correctlabel}[1]{\textcolor{okgreen}{\textbf{(#1)}\,\ding{51}}}
\newenvironment{qitem}[1]{\begin{list}{}{\setlength{\leftmargin}{3.3em}\setlength{\labelwidth}{3.1em}%
  \setlength{\labelsep}{0.2em}\setlength{\topsep}{3pt}\setlength{\parsep}{3pt}}\item[\textbf{Q.#1}\hfill]}{\end{list}}
\newcommand{\qsep}{\par\vspace{2pt}{\color{black!15}\hrule}\vspace{4pt}}
\newcommand{\qmeta}[2]{{\footnotesize\color{black!60}\textsc{#1} \;$\cdot$\; #2}\par}
\newtcolorbox{passagebox}[1]{enhanced, breakable, colback=brand!4, colframe=brand!70, fonttitle=\bfseries,
  title={#1}, left=6pt, right=6pt, top=4pt, bottom=4pt}
\newtcolorbox{sectionbanner}{enhanced, colback=brand, colframe=brand, coltext=white, fontupper=\large\bfseries,
  left=8pt, right=8pt, top=6pt, bottom=6pt, sharp corners}
\newtcolorbox{infobox}[1]{enhanced, breakable, colback=black!3, colframe=black!45, fonttitle=\bfseries,
  title={#1}, left=6pt, right=6pt, top=4pt, bottom=4pt}
\newtcolorbox{solbox}[1]{enhanced, breakable, colback=okgreen!3, colframe=okgreen!60!black, fonttitle=\bfseries,
  title={Answer: (#1) \quad Detailed solution}, left=6pt, right=6pt, top=3pt, bottom=4pt, before skip=4pt, after skip=6pt}
\newcommand{\VQEop}{\mathrm{VQE}}
\allowdisplaybreaks
\setlength{\emergencystretch}{3em}

\begin{document}
\pagestyle{fancy}\fancyhf{}
\fancyhead[L]{\small\color{black!70}Variational Quantum Eigensolver --- Assessment}
\fancyhead[R]{\small\color{black!70}\textbf{Question Paper}}
\fancyfoot[C]{\small Page \thepage\ of \pageref*{LastPage}}
\renewcommand{\headrulewidth}{0.4pt}

\thispagestyle{empty}
\begin{center}
{\color{brand}\rule{\linewidth}{1.2pt}}\\[6pt]
{\Huge\bfseries\color{brand} Variational Quantum Eigensolver}\\[4pt]
{\Large Assessment of the VQE algorithm}\\[4pt]
{\large\bfseries QUESTION PAPER}\\[2pt]
{\color{brand}\rule{\linewidth}{1.2pt}}\\[1pt]
\end{center}
\vspace{-2pt}
\begin{tabularx}{\linewidth}{@{}lXlr@{}}
\textbf{Time allowed:} & 3 hours & \textbf{Maximum marks:} & 240\\
\textbf{Questions:} & 80 (4 sections $\times$ 20) & \textbf{Answer key:} & separate Answer Book
\end{tabularx}
\vspace{6pt}

\textbf{Name of the candidate:} \hrulefill\quad \textbf{Roll No.:} \rule{35mm}{0.4pt}

\vspace{6pt}
\begin{infobox}{General instructions}
\begin{enumerate}[leftmargin=1.4em, itemsep=1pt, topsep=1pt]
\item This question paper has \textbf{four sections} of \textbf{20 questions} each:
\textbf{A}~Mathematical and Quantum Chemistry Foundations, \textbf{B}~Quantum Computing Foundations, \textbf{C}~VQE Fundamentals, \textbf{D}~Advanced VQE and Implementation.
\item Every question is of the \textbf{single-correct} type: four options, exactly one correct. Question formats used:
single-correct MCQ, matching lists (List-I/List-II), Assertion--Reason, Statement~I/II and paragraph-based questions.
\item For every question choose \textbf{ONE} option, (A), (B), (C) or (D), and record it against the question number in the answer script provided by the invigilator.
\item Calculators are \textbf{not} allowed. Rough work may be done on separate sheets.
\item Unless a question says otherwise, use the \textbf{conventions} printed below; they apply to every section.
\item The detailed solutions are in the separate \emph{Answer Book}; do not consult it before finishing the test.
\end{enumerate}
\end{infobox}
\begin{infobox}{Marking scheme (every question)}
Each question has \textbf{four} options (A), (B), (C) and (D). \textbf{ONLY ONE} of these four options is correct.\\[2pt]
\begin{tabular}{@{}ll@{}}
\emph{Full marks}: & $+3$ if ONLY the correct option is chosen;\\
\emph{Zero marks}: & $\phantom{+}0$ if none of the options is chosen (question unanswered);\\
\emph{Negative marks}: & $-1$ in all other cases (wrong option, or more than one option chosen).
\end{tabular}
\end{infobox}

\begin{infobox}{Conventions and notation used throughout the paper}
\small
\begin{itemize}[leftmargin=1.2em, itemsep=1pt, topsep=2pt]
\item $n$ qubits; a Pauli string is a word over $\{I, X, Y, Z\}$ whose character $k$ acts on qubit $k$; basis states $\ket{q_0q_1\dots q_{n-1}}$ are written
      \textbf{qubit 0 first}. $H = \sum_kc_kP_k$ with real $c_k$; $E_0$ = smallest eigenvalue; $E(\boldsymbol\theta) = \braket{\psi(\boldsymbol\theta)|H|\psi(\boldsymbol\theta)}$.
\item $R_Y(\theta) = e^{-i\theta Y/2}$, $R_X(\theta) = e^{-i\theta X/2}$, $R_Z(\theta) = e^{-i\theta Z/2}$, $S = \operatorname{diag}(1, i)$; CUDA-Q's \texttt{exp\_pauli(theta, q, word)} applies $e^{+i\theta P}$.
      Two-qubit matrices use the order $\ket{00}, \ket{01}, \ket{10}, \ket{11}$.
\item \textbf{Jordan--Wigner} (qubit $\ket1$ = occupied): $a_j^\dagger = Z_0\cdots Z_{j-1}\,\tfrac12(X_j - iY_j)$, $n_j = \tfrac12(I - Z_j)$.
\item \textbf{The notebook} (\texttt{VQE.ipynb}): default Ising chain $H = -J\sum_{i}Z_iZ_{i+1} - h\sum_iX_i$ (open, $N = 15$, $J = 1$, $h = 0.5$, 29 terms);
      hardware-efficient ansatz with $L = 2$ layers of [$R_Y$ on all qubits, CNOT ladder] plus a final $R_Y$ layer; COBYLA, 300 evaluations; exact
      energies via \texttt{cudaq.observe}. H$_2$ example: 4 qubits, 15 Pauli terms (STO-3G, Jordan--Wigner, spin-orbitals ordered $\sigma_g\alpha, \sigma_g\beta, \sigma_u\alpha, \sigma_u\beta$;
      coefficients correspond to $R = 0.740$ \AA), Hartree--Fock $\ket{1100}$, UCC state $e^{i\theta\,X_0X_1X_2Y_3}\ket{1100}$, FCI energy $-1.137284$ Ha.
      A \emph{previously recorded run} means the outputs stored in the notebook (CUDA-Q 0.16.0, CPU target \texttt{qpp-cpu}).
\item Energies of molecules in Hartree (Ha); 1 Ha $= 627.51$ kcal/mol; chemical accuracy $1.6$ mHa. Numerical answers are rounded.
\end{itemize}
\end{infobox}
\newpage

\begin{sectionbanner}SECTION A \enspace$\cdot$\enspace Mathematical and Quantum Chemistry Foundations \hfill Q.1--20\end{sectionbanner}

{\small \textbf{Syllabus:} eigenvalue problems, Rayleigh quotients, Hermitian matrices, the Schroedinger equation, Born--Oppenheimer, Slater determinants, second quantisation, Pauli decomposition, energy units and correlation energy.\\ \textbf{Formats:} Single-correct MCQ (Q.1--20); Matching lists (Q.3--11); Assertion--Reason (Q.7--18); Statement I / II (Q.8); Paragraph based (Q.14--15). Each question: $+3$ / $0$ / $-1$.}
\par\medskip

\needspace{7\baselineskip}
\begin{qitem}{1}
The ground-state energy (smallest eigenvalue) of $H = \begin{pmatrix}1 & 2\\ 2 & -2\end{pmatrix}$ and a normalised eigenvector for it are
\par\smallskip\noindent\begin{tabularx}{\linewidth}{@{}*{2}{L}@{}}
\optlabel{A}\,$-2$ and $(0, 1)^{\mathsf T}$ & \optlabel{B}\,$2$ and $\tfrac{1}{\sqrt5}(2, 1)^{\mathsf T}$ \\[3pt]
\optlabel{C}\,$-3$ and $\tfrac{1}{\sqrt5}(1, -2)^{\mathsf T}$ & \optlabel{D}\,$-3$ and $\tfrac{1}{\sqrt5}(2, 1)^{\mathsf T}$
\end{tabularx}
\end{qitem}
\qsep

\needspace{7\baselineskip}
\begin{qitem}{2}
For the same matrix $H = \begin{pmatrix}1 & 2\\ 2 & -2\end{pmatrix}$, the Rayleigh quotient $R(v) = \dfrac{v^{\mathsf T}Hv}{v^{\mathsf T}v}$ of the trial
vector $v = (1, -1)^{\mathsf T}$ is
\par\smallskip\noindent\begin{tabularx}{\linewidth}{@{}*{2}{L}@{}}
\optlabel{A}\,$-5$, which lies below $E_0$ & \optlabel{B}\,$-3$, exactly $E_0$ \\[3pt]
\optlabel{C}\,$-2.5$, which lies above $E_0 = -3$ & \optlabel{D}\,$0.5$
\end{tabularx}
\end{qitem}
\qsep

\needspace{7\baselineskip}
\begin{qitem}{3}
For the real symmetric matrix $M = \begin{pmatrix}2 & 1 & 0\\ 1 & 2 & 0\\ 0 & 0 & 5\end{pmatrix}$, match each quantity (List-I) with its value (List-II).
\par\smallskip\noindent{\renewcommand{\arraystretch}{1.35}%
\begin{tabularx}{\linewidth}{@{}c>{\raggedright\arraybackslash}p{0.33\linewidth}@{\hspace{1.2em}}cL@{}}
\toprule
\multicolumn{2}{@{}l}{\textbf{List-I}} & \multicolumn{2}{l@{}}{\textbf{List-II}} \\ \midrule
\textbf{(P)} & $\operatorname{tr} M$ & \textbf{(1)} & $15$ \\
\textbf{(Q)} & $\det M$ & \textbf{(2)} & $1$ \\
\textbf{(R)} & smallest eigenvalue of $M$ & \textbf{(3)} & $9$ \\
\textbf{(S)} & largest eigenvalue of $M$ & \textbf{(4)} & $5$ \\
\bottomrule
\end{tabularx}}
\par\smallskip\noindent\begin{tabularx}{\linewidth}{@{}*{2}{L}@{}}
\optlabel{A}\,P$\to$3, Q$\to$1, R$\to$4, S$\to$2 & \optlabel{B}\,P$\to$3, Q$\to$1, R$\to$2, S$\to$4 \\[3pt]
\optlabel{C}\,P$\to$1, Q$\to$3, R$\to$2, S$\to$4 & \optlabel{D}\,P$\to$3, Q$\to$4, R$\to$2, S$\to$1
\end{tabularx}
\end{qitem}
\qsep

\needspace{7\baselineskip}
\begin{qitem}{4}
An energy eigenstate $\ket{\psi_E}$ of a time-independent Hamiltonian ($H\ket{\psi_E} = E\ket{\psi_E}$) evolves under the time-dependent
Schr\"odinger equation $i\hbar\,\partial_t\ket{\psi(t)} = H\ket{\psi(t)}$. Which statement is correct?
\begin{itemize}[leftmargin=2.6em, labelsep=0.4em, itemsep=2pt, topsep=3pt]
\item[\optlabel{A}] $\ket{\psi(t)} = e^{-Et/\hbar}\ket{\psi_E}$: the state decays in time
\item[\optlabel{B}] The state is constant only if $E = 0$; otherwise it moves to other eigenstates
\item[\optlabel{C}] $\ket{\psi(t)} = e^{-iEt/\hbar}\ket{\psi_E}$: only a global phase changes, so all measurement probabilities are constant in time
\item[\optlabel{D}] The probabilities oscillate with angular frequency $E/\hbar$
\end{itemize}
\end{qitem}
\qsep

\needspace{7\baselineskip}
\begin{qitem}{5}
Using $1\ \text{Hartree} = 627.51$ kcal/mol, an energy error of $0.005$ Hartree corresponds to about ... and is ... chemically
accurate (chemical accuracy: 1 kcal/mol $\approx$ 1.6 mHa).
\par\smallskip\noindent\begin{tabularx}{\linewidth}{@{}*{2}{L}@{}}
\optlabel{A}\,$3.14$ kcal/mol; not & \optlabel{B}\,$3.14$ kcal/mol; still \\[3pt]
\optlabel{C}\,$0.31$ kcal/mol; still & \optlabel{D}\,$125.5$ kcal/mol; not
\end{tabularx}
\end{qitem}
\qsep

\needspace{7\baselineskip}
\begin{qitem}{6}
In the Born--Oppenheimer approximation the nuclei are held fixed and only the electrons are treated quantum mechanically. In the
notebook's 4-qubit H$_2$ Hamiltonian, the nuclear--nuclear repulsion energy appears
\begin{itemize}[leftmargin=2.6em, labelsep=0.4em, itemsep=2pt, topsep=3pt]
\item[\optlabel{A}] nowhere, because nuclear repulsion is neglected in quantum chemistry
\item[\optlabel{B}] spread over the four $\texttt{XXYY}$-type exchange terms
\item[\optlabel{C}] as the coefficient of the $\texttt{ZZII}$ term
\item[\optlabel{D}] inside the coefficient of the identity string $\texttt{IIII}$, as a constant energy shift
\end{itemize}
\end{qitem}
\qsep

\needspace{7\baselineskip}
\begin{qitem}{7}
\textbf{Assertion (A):} For every normalised state $\ket\psi$, $\braket{\psi|H|\psi} \ge E_0$, where $E_0$ is the smallest eigenvalue of the
Hermitian operator $H$.\\[1mm]
\textbf{Reason (R):} Expanding $\ket\psi = \sum_k c_k\ket{E_k}$ in the orthonormal eigenbasis of $H$ gives $\braket{\psi|H|\psi} = \sum_k |c_k|^2E_k$ with
$\sum_k|c_k|^2 = 1$.
\begin{itemize}[leftmargin=2.6em, labelsep=0.4em, itemsep=2pt, topsep=3pt]
\item[\optlabel{A}] Both Assertion (A) and Reason (R) are true, and R is the correct explanation of A
\item[\optlabel{B}] Both A and R are true, but R is \emph{not} the correct explanation of A
\item[\optlabel{C}] A is true, but R is false
\item[\optlabel{D}] A is false, but R is true
\end{itemize}
\end{qitem}
\qsep

\needspace{7\baselineskip}
\begin{qitem}{8}
\textbf{Statement I:} A real symmetric matrix can have complex (non-real) eigenvalues.\\[1mm]
\textbf{Statement II:} Eigenvectors of a Hermitian matrix that belong to different eigenvalues are orthogonal.
\begin{itemize}[leftmargin=2.6em, labelsep=0.4em, itemsep=2pt, topsep=3pt]
\item[\optlabel{A}] Both Statement I and Statement II are true
\item[\optlabel{B}] Both Statement I and Statement II are false
\item[\optlabel{C}] Statement I is true, but Statement II is false
\item[\optlabel{D}] Statement I is false, but Statement II is true
\end{itemize}
\end{qitem}
\qsep

\needspace{7\baselineskip}
\begin{qitem}{9}
H$_2$ in the STO-3G basis has 2 spatial orbitals, i.e.\ 4 spin-orbitals, and 2 electrons. The number of 2-electron Slater determinants,
and the number of them with one spin-up and one spin-down electron ($S_z = 0$), are
\par\smallskip\noindent\begin{tabularx}{\linewidth}{@{}*{4}{L}@{}}
\optlabel{A}\,$6$ and $2$ & \optlabel{B}\,$6$ and $4$ & \optlabel{C}\,$16$ and $4$ & \optlabel{D}\,$4$ and $2$
\end{tabularx}
\end{qitem}
\qsep

\needspace{7\baselineskip}
\begin{qitem}{10}
LiH in the minimal STO-3G basis has one $1s$ orbital on H and the orbitals $1s, 2s, 2p_x, 2p_y, 2p_z$ on Li. With the Jordan--Wigner
mapping (one qubit per spin-orbital, no reductions), the number of qubits and the number of electrons are
\par\smallskip\noindent\begin{tabularx}{\linewidth}{@{}*{2}{L}@{}}
\optlabel{A}\,$24$ qubits, $4$ electrons & \optlabel{B}\,$12$ qubits, $3$ electrons \\[3pt]
\optlabel{C}\,$12$ qubits, $4$ electrons & \optlabel{D}\,$6$ qubits, $4$ electrons
\end{tabularx}
\end{qitem}
\qsep

\needspace{7\baselineskip}
\begin{qitem}{11}
Fermionic creation and annihilation operators $a_p^\dagger$, $a_p$ satisfy the canonical anticommutation relations. Match each expression
(List-I) with its value (List-II); $\ket{\mathrm{vac}}$ is the empty (vacuum) state.
\par\smallskip\noindent{\renewcommand{\arraystretch}{1.35}%
\begin{tabularx}{\linewidth}{@{}c>{\raggedright\arraybackslash}p{0.33\linewidth}@{\hspace{1.2em}}cL@{}}
\toprule
\multicolumn{2}{@{}l}{\textbf{List-I}} & \multicolumn{2}{l@{}}{\textbf{List-II}} \\ \midrule
\textbf{(P)} & $\{a_p, a_q^\dagger\} = a_pa_q^\dagger + a_q^\dagger a_p$ & \textbf{(1)} & $\{0, 1\}$ \\
\textbf{(Q)} & $a_p^\dagger a_p^\dagger$ & \textbf{(2)} & $\delta_{pq}\,I$ \\
\textbf{(R)} & spectrum of $n_p = a_p^\dagger a_p$ & \textbf{(3)} & $0$ (the zero vector) \\
\textbf{(S)} & $a_p\ket{\mathrm{vac}}$ & \textbf{(4)} & $0$ (the zero operator) \\
\bottomrule
\end{tabularx}}
\par\smallskip\noindent\begin{tabularx}{\linewidth}{@{}*{2}{L}@{}}
\optlabel{A}\,P$\to$2, Q$\to$4, R$\to$1, S$\to$3 & \optlabel{B}\,P$\to$2, Q$\to$3, R$\to$1, S$\to$4 \\[3pt]
\optlabel{C}\,P$\to$2, Q$\to$4, R$\to$3, S$\to$1 & \optlabel{D}\,P$\to$4, Q$\to$2, R$\to$1, S$\to$3
\end{tabularx}
\end{qitem}
\qsep

\needspace{7\baselineskip}
\begin{qitem}{12}
In second quantisation $H = \sum_{pq}h_{pq}a_p^\dagger a_q + \tfrac12\sum_{pqrs}h_{pqrs}a_p^\dagger a_q^\dagger a_ra_s$ (+ constant). For $N$ spin-orbitals the number of
terms in this Hamiltonian grows as
\begin{itemize}[leftmargin=2.6em, labelsep=0.4em, itemsep=2pt, topsep=3pt]
\item[\optlabel{A}] $O(N)$, one term per orbital
\item[\optlabel{B}] $O(N^4)$, from the two-electron integrals $h_{pqrs}$
\item[\optlabel{C}] $O(N^2)$, from the one-electron integrals only
\item[\optlabel{D}] $O(2^N)$, like the dimension of the Fock space
\end{itemize}
\end{qitem}
\qsep

\needspace{7\baselineskip}
\begin{qitem}{13}
For a Pauli operator $P$ ($P^2 = I$), $e^{i\theta P} = \cos\theta\,I + i\sin\theta\,P$. Hence $e^{i\pi X/4}$ equals
\par\smallskip\noindent\begin{tabularx}{\linewidth}{@{}*{4}{L}@{}}
\optlabel{A}\,$\begin{pmatrix}0 & i\\ i & 0\end{pmatrix}$ & \optlabel{B}\,$\dfrac{1}{\sqrt2}\begin{pmatrix}1 & 1\\ 1 & -1\end{pmatrix}$ & \optlabel{C}\,$\dfrac{1}{\sqrt2}\begin{pmatrix}1 & i\\ i & 1\end{pmatrix}$ & \optlabel{D}\,$\dfrac{1}{\sqrt2}\begin{pmatrix}1 & -i\\ -i & 1\end{pmatrix}$
\end{tabularx}
\end{qitem}
\qsep

\needspace{20\baselineskip}
\begin{passagebox}{Paragraph for Questions 14 and 15}
A single qubit has the Hamiltonian $H = -Z + \tfrac12 X = \begin{pmatrix}-1 & \tfrac12\\ \tfrac12 & 1\end{pmatrix}$.
A one-parameter trial state is prepared as $\ket{\psi(\theta)} = R_Y(\theta)\ket0 = \cos\tfrac\theta2\ket0 + \sin\tfrac\theta2\ket1$, and its
energy $E(\theta) = \braket{\psi(\theta)|H|\psi(\theta)}$ is minimised over $\theta \in (-\pi, \pi]$.
\end{passagebox}
\needspace{7\baselineskip}
\begin{qitem}{14}
The energy function $E(\theta)$ and the exact ground-state energy $E_0$ are
\begin{itemize}[leftmargin=2.6em, labelsep=0.4em, itemsep=2pt, topsep=3pt]
\item[\optlabel{A}] $E(\theta) = -\cos\theta + \tfrac12\sin\theta$ and $E_0 = -1$
\item[\optlabel{B}] $E(\theta) = -\cos\tfrac\theta2 + \tfrac12\sin\tfrac\theta2$ and $E_0 = -\tfrac{\sqrt5}{2}$
\item[\optlabel{C}] $E(\theta) = -\cos\theta + \sin\theta$ and $E_0 = -\sqrt2$
\item[\optlabel{D}] $E(\theta) = -\cos\theta + \tfrac12\sin\theta$ and $E_0 = -\tfrac{\sqrt5}{2} \approx -1.118$
\end{itemize}
\end{qitem}
\qsep

\needspace{7\baselineskip}
\begin{qitem}{15}
The optimal angle $\theta^*$, and the energy error $E(0) - E_0$ of the starting point $\theta = 0$ (the state $\ket0$), are
\par\smallskip\noindent\begin{tabularx}{\linewidth}{@{}*{2}{L}@{}}
\optlabel{A}\,$\theta^* = -\arctan\tfrac12 \approx -0.464$ rad; error $0$ & \optlabel{B}\,$\theta^* = -\arctan 2 \approx -1.107$ rad; error $\approx 0.118$ \\[3pt]
\optlabel{C}\,$\theta^* = +\arctan\tfrac12 \approx 0.464$ rad; error $\approx 0.118$ & \optlabel{D}\,$\theta^* = -\arctan\tfrac12 \approx -0.464$ rad; error $\approx 0.118$
\end{tabularx}
\end{qitem}
\qsep

\needspace{7\baselineskip}
\begin{qitem}{16}
The Hermitian matrix $M = \begin{pmatrix}3 & 1 - 2i\\ 1 + 2i & -1\end{pmatrix}$ written as $aI + bX + cY + dZ$ is
\par\smallskip\noindent\begin{tabularx}{\linewidth}{@{}*{4}{L}@{}}
\optlabel{A}\,$I + X + 2Y + 2Z$ & \optlabel{B}\,$3I + X + 2Y - Z$ & \optlabel{C}\,$I + X - 2Y + 2Z$ & \optlabel{D}\,$2I + X + 2Y + Z$
\end{tabularx}
\end{qitem}
\qsep

\needspace{7\baselineskip}
\begin{qitem}{17}
The diagonal two-qubit Hamiltonian $H = \operatorname{diag}(1, 2, 3, 4)$ (basis $\ket{00}, \ket{01}, \ket{10}, \ket{11}$) written as a sum of Pauli strings is
\par\smallskip\noindent\begin{tabularx}{\linewidth}{@{}*{2}{L}@{}}
\optlabel{A}\,$2.5\,II - 0.5\,ZI - 1.0\,IZ$ & \optlabel{B}\,$2.5\,II - 1.0\,ZI - 0.5\,IZ + 0.5\,ZZ$ \\[3pt]
\optlabel{C}\,$2.5\,II + 1.0\,ZI + 0.5\,IZ$ & \optlabel{D}\,$2.5\,II - 1.0\,ZI - 0.5\,IZ$
\end{tabularx}
\end{qitem}
\qsep

\needspace{7\baselineskip}
\begin{qitem}{18}
\textbf{Assertion (A):} For H$_2$ at large bond length, the restricted Hartree--Fock (single closed-shell determinant) energy lies far above
the exact energy.\\[1mm]
\textbf{Reason (R):} The nuclear--nuclear repulsion $1/R$ tends to zero as $R \to \infty$.
\begin{itemize}[leftmargin=2.6em, labelsep=0.4em, itemsep=2pt, topsep=3pt]
\item[\optlabel{A}] Both Assertion (A) and Reason (R) are true, and R is the correct explanation of A
\item[\optlabel{B}] Both A and R are true, but R is \emph{not} the correct explanation of A
\item[\optlabel{C}] A is true, but R is false
\item[\optlabel{D}] A is false, but R is true
\end{itemize}
\end{qitem}
\qsep

\needspace{7\baselineskip}
\begin{qitem}{19}
For the notebook's H$_2$ Hamiltonian, the Hartree--Fock energy is $E_{\mathrm{HF}} = \braket{1100|H|1100} = -1.116759$ Ha and the FCI energy is
$-1.137284$ Ha. The correlation energy $E_{\mathrm{corr}} = E_{\mathrm{FCI}} - E_{\mathrm{HF}}$ is about
\par\smallskip\noindent\begin{tabularx}{\linewidth}{@{}*{4}{L}@{}}
\optlabel{A}\,$-20.5$ mHa & \optlabel{B}\,$+20.5$ mHa & \optlabel{C}\,$-2.05$ mHa & \optlabel{D}\,$-1137$ mHa
\end{tabularx}
\end{qitem}
\qsep

\needspace{7\baselineskip}
\begin{qitem}{20}
A trial state of H$_2$ has $90\,\%$ weight on the ground state ($E_0 = -1.13728$ Ha) and $10\,\%$ on another eigenstate of the 4-qubit Hamiltonian with energy
$E_1 = -0.53821$ Ha. Its energy expectation and its error above $E_0$ are
\par\smallskip\noindent\begin{tabularx}{\linewidth}{@{}*{2}{L}@{}}
\optlabel{A}\,$-1.1971$ Ha and $-59.9$ mHa & \optlabel{B}\,$-1.0774$ Ha and $59.9$ mHa \\[3pt]
\optlabel{C}\,$-0.8378$ Ha and $299.5$ mHa & \optlabel{D}\,$-1.1373$ Ha and $0$ mHa
\end{tabularx}
\end{qitem}
\qsep

\newpage

\begin{sectionbanner}SECTION B \enspace$\cdot$\enspace Quantum Computing Foundations \hfill Q.21--40\end{sectionbanner}

{\small \textbf{Syllabus:} qubits and the Bloch sphere, rotation gates, Pauli algebra and tensor products, measurement bases, CNOT and entanglement, qubit ordering, \texttt{exp\_pauli}, commuting observables, memory and variance.\\ \textbf{Formats:} Single-correct MCQ (Q.21--40); Matching lists (Q.23--24); Assertion--Reason (Q.32); Statement I / II (Q.28--39); Paragraph based (Q.36--37). Each question: $+3$ / $0$ / $-1$.}
\par\medskip

\needspace{7\baselineskip}
\begin{qitem}{21}
The state $R_Y(\pi/3)\ket0$ is measured in the computational basis. The probability of the outcome 1 is
\par\smallskip\noindent\begin{tabularx}{\linewidth}{@{}*{4}{L}@{}}
\optlabel{A}\,$\tfrac34$ & \optlabel{B}\,$\tfrac{\sqrt3}{2}$ & \optlabel{C}\,$\tfrac12$ & \optlabel{D}\,$\tfrac14$
\end{tabularx}
\end{qitem}
\qsep

\needspace{7\baselineskip}
\begin{qitem}{22}
The two-qubit operator $(X\otimes Y)(Y\otimes X)$ equals
\par\smallskip\noindent\begin{tabularx}{\linewidth}{@{}*{4}{L}@{}}
\optlabel{A}\,$Z\otimes Z$ & \optlabel{B}\,$Y\otimes Y$ & \optlabel{C}\,$-Z\otimes Z$ & \optlabel{D}\,$i\,Z\otimes Z$
\end{tabularx}
\end{qitem}
\qsep

\needspace{7\baselineskip}
\begin{qitem}{23}
Match each single-qubit state (List-I) with its Bloch vector $(\braket X, \braket Y, \braket Z)$ (List-II).
\par\smallskip\noindent{\renewcommand{\arraystretch}{1.35}%
\begin{tabularx}{\linewidth}{@{}c>{\raggedright\arraybackslash}p{0.33\linewidth}@{\hspace{1.2em}}cL@{}}
\toprule
\multicolumn{2}{@{}l}{\textbf{List-I}} & \multicolumn{2}{l@{}}{\textbf{List-II}} \\ \midrule
\textbf{(P)} & $\ket0$ & \textbf{(1)} & $(0, 1, 0)$ \\
\textbf{(Q)} & $\ket- = \tfrac{1}{\sqrt2}(\ket0 - \ket1)$ & \textbf{(2)} & $(\tfrac{\sqrt3}{2}, 0, \tfrac12)$ \\
\textbf{(R)} & $\tfrac{1}{\sqrt2}(\ket0 + i\ket1)$ & \textbf{(3)} & $(0, 0, 1)$ \\
\textbf{(S)} & $R_Y(\pi/3)\ket0$ & \textbf{(4)} & $(-1, 0, 0)$ \\
\bottomrule
\end{tabularx}}
\par\smallskip\noindent\begin{tabularx}{\linewidth}{@{}*{2}{L}@{}}
\optlabel{A}\,P$\to$3, Q$\to$4, R$\to$2, S$\to$1 & \optlabel{B}\,P$\to$3, Q$\to$1, R$\to$4, S$\to$2 \\[3pt]
\optlabel{C}\,P$\to$4, Q$\to$3, R$\to$1, S$\to$2 & \optlabel{D}\,P$\to$3, Q$\to$4, R$\to$1, S$\to$2
\end{tabularx}
\end{qitem}
\qsep

\needspace{7\baselineskip}
\begin{qitem}{24}
To estimate $\braket P$ from computational-basis measurements, each Pauli string is first rotated into a string of $Z$'s and $I$'s. Match each
observable (List-I) with the gates applied before measuring (List-II).
\par\smallskip\noindent{\renewcommand{\arraystretch}{1.35}%
\begin{tabularx}{\linewidth}{@{}c>{\raggedright\arraybackslash}p{0.33\linewidth}@{\hspace{1.2em}}cL@{}}
\toprule
\multicolumn{2}{@{}l}{\textbf{List-I}} & \multicolumn{2}{l@{}}{\textbf{List-II}} \\ \midrule
\textbf{(P)} & $X$ & \textbf{(1)} & $H$ on qubit 0 only \\
\textbf{(Q)} & $Y$ & \textbf{(2)} & nothing \\
\textbf{(R)} & $Z$ & \textbf{(3)} & $H$ \\
\textbf{(S)} & $X_0Z_1$ & \textbf{(4)} & $S^\dagger$, then $H$ \\
\bottomrule
\end{tabularx}}
\par\smallskip\noindent\begin{tabularx}{\linewidth}{@{}*{2}{L}@{}}
\optlabel{A}\,P$\to$3, Q$\to$2, R$\to$4, S$\to$1 & \optlabel{B}\,P$\to$3, Q$\to$4, R$\to$1, S$\to$2 \\[3pt]
\optlabel{C}\,P$\to$3, Q$\to$4, R$\to$2, S$\to$1 & \optlabel{D}\,P$\to$4, Q$\to$3, R$\to$2, S$\to$1
\end{tabularx}
\end{qitem}
\qsep

\needspace{7\baselineskip}
\begin{qitem}{25}
For the Bell state $\ket{\Phi^+} = \tfrac{1}{\sqrt2}(\ket{00} + \ket{11})$, the expectation value $\braket{Y\otimes Y}$ is
\par\smallskip\noindent\begin{tabularx}{\linewidth}{@{}*{4}{L}@{}}
\optlabel{A}\,$+1$ & \optlabel{B}\,$-\tfrac12$ & \optlabel{C}\,$0$ & \optlabel{D}\,$-1$
\end{tabularx}
\end{qitem}
\qsep

\needspace{7\baselineskip}
\begin{qitem}{26}
The Hartree--Fock state of the notebook's H$_2$ example is $\ket{1100}$, written with qubit 0 first. In the vector returned by \texttt{cudaq.get\_state},
which lists amplitudes with qubit 0 as the \emph{least} significant bit, this basis state has the index
\par\smallskip\noindent\begin{tabularx}{\linewidth}{@{}*{4}{L}@{}}
\optlabel{A}\,$12$ & \optlabel{B}\,$3$ & \optlabel{C}\,$0$ & \optlabel{D}\,$6$
\end{tabularx}
\end{qitem}
\qsep

\needspace{7\baselineskip}
\begin{qitem}{27}
Starting from $\ket{00}$, apply $R_Y(\pi/2)$ to qubit 0 and then a CNOT with control 0 and target 1. The resulting state is
\par\smallskip\noindent\begin{tabularx}{\linewidth}{@{}*{2}{L}@{}}
\optlabel{A}\,$\tfrac{1}{\sqrt2}(\ket{00} + \ket{11})$ & \optlabel{B}\,$\tfrac{1}{\sqrt2}(\ket{00} + \ket{10})$ \\[3pt]
\optlabel{C}\,$\tfrac{1}{\sqrt2}(\ket{00} - \ket{11})$ & \optlabel{D}\,$\tfrac12(\ket{00} + \ket{01} + \ket{10} + \ket{11})$
\end{tabularx}
\end{qitem}
\qsep

\needspace{7\baselineskip}
\begin{qitem}{28}
\textbf{Statement I:} Any circuit of gates applied to $\ket{0\dots0}$ prepares a normalised state, so the ansatz state $\ket{\psi(\boldsymbol\theta)}$ needs no
explicit normalisation.\\[1mm]
\textbf{Statement II:} Measuring in the computational basis directly reveals the eigenvalue of an $X_i$ term, without any basis change.
\begin{itemize}[leftmargin=2.6em, labelsep=0.4em, itemsep=2pt, topsep=3pt]
\item[\optlabel{A}] Both Statement I and Statement II are true
\item[\optlabel{B}] Both Statement I and Statement II are false
\item[\optlabel{C}] Statement I is true, but Statement II is false
\item[\optlabel{D}] Statement I is false, but Statement II is true
\end{itemize}
\end{qitem}
\qsep

\needspace{7\baselineskip}
\begin{qitem}{29}
Consider the four exchange terms of the notebook's H$_2$ Hamiltonian, $\texttt{XXYY}$, $\texttt{XYYX}$, $\texttt{YXXY}$, $\texttt{YYXX}$. Which statement is
correct?
\begin{itemize}[leftmargin=2.6em, labelsep=0.4em, itemsep=2pt, topsep=3pt]
\item[\optlabel{A}] Every pair anticommutes
\item[\optlabel{B}] All six pairs commute, but no pair is qubit-wise commuting
\item[\optlabel{C}] Only $\texttt{XXYY}$ and $\texttt{YYXX}$ commute
\item[\optlabel{D}] All six pairs are qubit-wise commuting, so one measurement setting suffices
\end{itemize}
\end{qitem}
\qsep

\needspace{7\baselineskip}
\begin{qitem}{30}
CUDA-Q's \texttt{exp\_pauli(theta, q, word)} applies $e^{i\theta P}$ for the Pauli string $P$ given by \texttt{word}. Applied with the word
$\texttt{"ZZ"}$ to the basis state $\ket{01}$, it gives
\par\smallskip\noindent\begin{tabularx}{\linewidth}{@{}*{2}{L}@{}}
\optlabel{A}\,$\cos\theta\ket{01} + i\sin\theta\ket{10}$ & \optlabel{B}\,$e^{+i\theta}\ket{01}$ \\[3pt]
\optlabel{C}\,$e^{-i\theta}\ket{01}$ & \optlabel{D}\,$\ket{10}$
\end{tabularx}
\end{qitem}
\qsep

\needspace{7\baselineskip}
\begin{qitem}{31}
The four-qubit Pauli string $X_0X_1X_2Y_3$ (word $\texttt{"XXXY"}$) acting on $\ket{1100}$ gives
\par\smallskip\noindent\begin{tabularx}{\linewidth}{@{}*{4}{L}@{}}
\optlabel{A}\,$i\ket{0011}$ & \optlabel{B}\,$\ket{0011}$ & \optlabel{C}\,$-i\ket{0011}$ & \optlabel{D}\,$i\ket{1100}$
\end{tabularx}
\end{qitem}
\qsep

\needspace{7\baselineskip}
\begin{qitem}{32}
\textbf{Assertion (A):} For a Hermitian $H$, the expectation value $\braket{\psi|H|\psi}$ is a real number for every state $\ket\psi$.\\[1mm]
\textbf{Reason (R):} The identity $\braket{\psi|H|\psi} = \braket{\psi|H|\psi}^*$ holds only for normalised states.
\begin{itemize}[leftmargin=2.6em, labelsep=0.4em, itemsep=2pt, topsep=3pt]
\item[\optlabel{A}] Both Assertion (A) and Reason (R) are true, and R is the correct explanation of A
\item[\optlabel{B}] Both A and R are true, but R is \emph{not} the correct explanation of A
\item[\optlabel{C}] A is true, but R is false
\item[\optlabel{D}] A is false, but R is true
\end{itemize}
\end{qitem}
\qsep

\needspace{7\baselineskip}
\begin{qitem}{33}
A complex fp64 amplitude needs 16 bytes. For LiH in STO-3G with Jordan--Wigner (12 qubits), the memory needed for the state vector and
for the density matrix ($2^n\times2^n$ amplitudes) is
\par\smallskip\noindent\begin{tabularx}{\linewidth}{@{}*{2}{L}@{}}
\optlabel{A}\,$192$ B and $2.3$ kB & \optlabel{B}\,$65.5$ kB and $131$ kB \\[3pt]
\optlabel{C}\,$268$ MB and $65.5$ kB & \optlabel{D}\,$65.5$ kB and $268$ MB
\end{tabularx}
\end{qitem}
\qsep

\needspace{7\baselineskip}
\begin{qitem}{34}
Two qubits start in $\ket{00}$; $R_Y(\theta_0)$ and $R_Y(\theta_1)$ are applied to qubits 0 and 1, followed by a CNOT (control 0, target 1). The
expectation value $\braket{Z_0Z_1}$ of the final state is
\par\smallskip\noindent\begin{tabularx}{\linewidth}{@{}*{4}{L}@{}}
\optlabel{A}\,$\cos(\theta_0 + \theta_1)$ & \optlabel{B}\,$\cos\theta_1$ & \optlabel{C}\,$\cos\theta_0$ & \optlabel{D}\,$\cos\theta_0\cos\theta_1$
\end{tabularx}
\end{qitem}
\qsep

\needspace{7\baselineskip}
\begin{qitem}{35}
Which unitary $U$ satisfies $UYU^\dagger = Z$, so that measuring $Z$ after applying $U$ measures $Y$?
\par\smallskip\noindent\begin{tabularx}{\linewidth}{@{}*{2}{L}@{}}
\optlabel{A}\,$U = S$ & \optlabel{B}\,$U = HS^\dagger$ (first $S^\dagger$, then $H$) \\[3pt]
\optlabel{C}\,$U = SH$ (first $H$, then $S$) & \optlabel{D}\,$U = H$
\end{tabularx}
\end{qitem}
\qsep

\needspace{20\baselineskip}
\begin{passagebox}{Paragraph for Questions 36 and 37}
A two-qubit transverse-field Ising Hamiltonian $H = -Z_0Z_1 - \tfrac12(X_0 + X_1)$ (the notebook's \texttt{"tfim"} model with $N = 2$,
$J = 1$, $h = 0.5$) is treated with a tiny hardware-efficient circuit: starting from $\ket{00}$, apply $R_Y(\theta_0)$ to qubit 0 and $R_Y(\theta_1)$
to qubit 1, then a CNOT with control 0 and target 1. The angles are $\theta_0 = \pi/2$, $\theta_1 = 0$.
\end{passagebox}
\needspace{7\baselineskip}
\begin{qitem}{36}
The energy $\braket{\psi|H|\psi}$ of the prepared state is
\par\smallskip\noindent\begin{tabularx}{\linewidth}{@{}*{4}{L}@{}}
\optlabel{A}\,$-1$ & \optlabel{B}\,$0$ & \optlabel{C}\,$-\sqrt2 \approx -1.414$ & \optlabel{D}\,$-1.5$
\end{tabularx}
\end{qitem}
\qsep

\needspace{7\baselineskip}
\begin{qitem}{37}
The exact ground-state energy $E_0$ of this $H$, and the relative error of the prepared state, are
\par\smallskip\noindent\begin{tabularx}{\linewidth}{@{}*{2}{L}@{}}
\optlabel{A}\,$-1$; $0\,\%$ & \optlabel{B}\,$-\sqrt2 \approx -1.414$; $29.3\,\%$ \\[3pt]
\optlabel{C}\,$-\sqrt2$; $41.4\,\%$ & \optlabel{D}\,$-1.5$; $33.3\,\%$
\end{tabularx}
\end{qitem}
\qsep

\needspace{7\baselineskip}
\begin{qitem}{38}
Every Hermitian $2^n\times2^n$ matrix is a real linear combination of the $n$-qubit Pauli strings. For $n = 4$ the number of Pauli strings,
and the number actually needed for the notebook's H$_2$ Hamiltonian, are
\par\smallskip\noindent\begin{tabularx}{\linewidth}{@{}*{4}{L}@{}}
\optlabel{A}\,$256$ and $256$ & \optlabel{B}\,$64$ and $15$ & \optlabel{C}\,$16$ and $15$ & \optlabel{D}\,$256$ and $15$
\end{tabularx}
\end{qitem}
\qsep

\needspace{7\baselineskip}
\begin{qitem}{39}
\textbf{Statement I:} Two Pauli strings either commute or anticommute.\\[1mm]
\textbf{Statement II:} A set of mutually commuting Pauli strings can be measured simultaneously after one common basis change.
\begin{itemize}[leftmargin=2.6em, labelsep=0.4em, itemsep=2pt, topsep=3pt]
\item[\optlabel{A}] Both Statement I and Statement II are true
\item[\optlabel{B}] Both Statement I and Statement II are false
\item[\optlabel{C}] Statement I is true, but Statement II is false
\item[\optlabel{D}] Statement I is false, but Statement II is true
\end{itemize}
\end{qitem}
\qsep

\needspace{7\baselineskip}
\begin{qitem}{40}
For $H = -Z + \tfrac12X$ (one qubit), the energy variance $\mathrm{Var}(H) = \braket{H^2} - \braket H^2$ in the state $\ket0$, and in the exact ground state, are
\par\smallskip\noindent\begin{tabularx}{\linewidth}{@{}*{4}{L}@{}}
\optlabel{A}\,$0.25$ and $1.25$ & \optlabel{B}\,$1.25$ and $0$ & \optlabel{C}\,$0.25$ and $0$ & \optlabel{D}\,$0$ and $0.25$
\end{tabularx}
\end{qitem}
\qsep

\newpage

\begin{sectionbanner}SECTION C \enspace$\cdot$\enspace VQE Fundamentals \hfill Q.41--60\end{sectionbanner}

{\small \textbf{Syllabus:} variational and Rayleigh--Ritz principles, ansatz design, Hartree--Fock reference, expectation values, Hamiltonian decomposition, measurement grouping, classical optimizers, convergence and the notebook's code.\\ \textbf{Formats:} Single-correct MCQ (Q.41--57); Matching lists (Q.46--58); Assertion--Reason (Q.49); Statement I / II (Q.52); Paragraph based (Q.59--60). Each question: $+3$ / $0$ / $-1$.}
\par\medskip

\needspace{7\baselineskip}
\begin{qitem}{41}
A VQE run with \emph{exact} expectation values (no shots, no noise) reports $E(\boldsymbol\theta) = -15.10$ for the 15-qubit Ising chain whose exact ground
energy is $E_0 = -15.0825$. The most likely explanation is
\begin{itemize}[leftmargin=2.6em, labelsep=0.4em, itemsep=2pt, topsep=3pt]
\item[\optlabel{A}] the optimizer found a state better than the exact ground state
\item[\optlabel{B}] a deeper ansatz always produces energies below the exact value
\item[\optlabel{C}] an error in the Hamiltonian, the reference energy or the state normalisation --- a correct calculation cannot go below $E_0$
\item[\optlabel{D}] shot noise pushed the estimate below $E_0$
\end{itemize}
\end{qitem}
\qsep

\needspace{7\baselineskip}
\begin{qitem}{42}
In the Rayleigh--Ritz method, $H$ is diagonalised inside an $m$-dimensional trial subspace, giving Ritz values $\varepsilon_1 \le \dots \le \varepsilon_m$. Which
statement is correct?
\begin{itemize}[leftmargin=2.6em, labelsep=0.4em, itemsep=2pt, topsep=3pt]
\item[\optlabel{A}] Only $\varepsilon_1$ is bounded; the other Ritz values can lie below the exact eigenvalues
\item[\optlabel{B}] The Ritz values are lower bounds to the exact eigenvalues
\item[\optlabel{C}] $\varepsilon_k \ge E_{k-1}$ for every $k$: each Ritz value is an upper bound to the corresponding exact eigenvalue
\item[\optlabel{D}] The Ritz values are exact whenever $m \ge 2$
\end{itemize}
\end{qitem}
\qsep

\needspace{7\baselineskip}
\begin{qitem}{43}
The notebook's hardware-efficient ansatz uses $L$ layers of [$R_Y$ on every qubit, CNOT ladder over neighbours] plus a final $R_Y$ layer. For
$N = 10$ qubits and $L = 3$, the numbers of parameters, CNOTs and gates are
\par\smallskip\noindent\begin{tabularx}{\linewidth}{@{}*{4}{L}@{}}
\optlabel{A}\,$40$, $30$ and $70$ & \optlabel{B}\,$30$, $27$ and $57$ & \optlabel{C}\,$40$, $9$ and $49$ & \optlabel{D}\,$40$, $27$ and $67$
\end{tabularx}
\end{qitem}
\qsep

\needspace{7\baselineskip}
\begin{qitem}{44}
The UCC kernel starts with \texttt{x(qubits[0]); x(qubits[1])}. In the Jordan--Wigner picture with spin-orbitals ordered (orbital 0 $\alpha$, orbital 0 $\beta$,
orbital 1 $\alpha$, orbital 1 $\beta$), this prepares
\begin{itemize}[leftmargin=2.6em, labelsep=0.4em, itemsep=2pt, topsep=3pt]
\item[\optlabel{A}] the state with both electrons in the antibonding orbital
\item[\optlabel{B}] an equal superposition of all 2-electron determinants
\item[\optlabel{C}] the Hartree--Fock determinant $\ket{1100}$: both electrons in the lower (bonding) spatial orbital, one with each spin
\item[\optlabel{D}] a triplet state with both electrons spin-up
\end{itemize}
\end{qitem}
\qsep

\needspace{7\baselineskip}
\begin{qitem}{45}
Let $\ket u$, $\ket v$ be orthonormal, $H_{uu} = a$, $H_{vv} = c$, $H_{uv} = H_{vu} = b$ (real). For the trial state $\ket{\psi(\theta)} = \cos\theta\ket u - \sin\theta\ket v$, the energy is
\par\smallskip\noindent\begin{tabularx}{\linewidth}{@{}*{2}{L}@{}}
\optlabel{A}\,$a\cos^2\theta + c\sin^2\theta - b\sin2\theta$ & \optlabel{B}\,$a\cos^2\theta - c\sin^2\theta$ \\[3pt]
\optlabel{C}\,$a\cos\theta + c\sin\theta - b$ & \optlabel{D}\,$a\cos^2\theta + c\sin^2\theta + b\sin2\theta$
\end{tabularx}
\end{qitem}
\qsep

\needspace{7\baselineskip}
\begin{qitem}{46}
Match each VQE ingredient (List-I) with the code in \texttt{VQE.ipynb} that implements it (List-II).
\par\smallskip\noindent{\renewcommand{\arraystretch}{1.35}%
\begin{tabularx}{\linewidth}{@{}c>{\raggedright\arraybackslash}p{0.33\linewidth}@{\hspace{1.2em}}cL@{}}
\toprule
\multicolumn{2}{@{}l}{\textbf{List-I}} & \multicolumn{2}{l@{}}{\textbf{List-II}} \\ \midrule
\textbf{(P)} & Hartree--Fock reference state & \textbf{(1)} & \texttt{minimize(energy, initial\_theta, method="COBYLA", ...)} \\
\textbf{(Q)} & energy $E(\boldsymbol\theta) = \braket{\psi(\boldsymbol\theta)|H|\psi(\boldsymbol\theta)}$ & \textbf{(2)} & \texttt{coefficient * cudaq.SpinOperator.from\_word(word)}, summed \\
\textbf{(R)} & classical parameter update & \textbf{(3)} & \texttt{x(qubits[0]); x(qubits[1])} \\
\textbf{(S)} & Hamiltonian as a sum of Pauli strings & \textbf{(4)} & \texttt{cudaq.observe(kernel, hamiltonian, *args).expectation()} \\
\bottomrule
\end{tabularx}}
\par\smallskip\noindent\begin{tabularx}{\linewidth}{@{}*{2}{L}@{}}
\optlabel{A}\,P$\to$3, Q$\to$4, R$\to$1, S$\to$2 & \optlabel{B}\,P$\to$3, Q$\to$4, R$\to$2, S$\to$1 \\[3pt]
\optlabel{C}\,P$\to$4, Q$\to$3, R$\to$1, S$\to$2 & \optlabel{D}\,P$\to$3, Q$\to$1, R$\to$4, S$\to$2
\end{tabularx}
\end{qitem}
\qsep

\needspace{7\baselineskip}
\begin{qitem}{47}
$H = 0.5\,Z_0 - 0.3\,Z_1 + 0.2\,X_0X_1$. A trial state has $\braket{Z_0} = 0.6$, $\braket{Z_1} = -0.2$ and $\braket{X_0X_1} = 0.5$. Its energy is
\par\smallskip\noindent\begin{tabularx}{\linewidth}{@{}*{4}{L}@{}}
\optlabel{A}\,$0.34$ & \optlabel{B}\,$1.3$ & \optlabel{C}\,$0.46$ & \optlabel{D}\,$0.26$
\end{tabularx}
\end{qitem}
\qsep

\needspace{7\baselineskip}
\begin{qitem}{48}
Only the identity and the $Z$-type terms of the notebook's H$_2$ Hamiltonian contribute to $\braket{1100|H|1100}$. With coefficients
$c_{IIII} = -0.097066$, $c_{ZIII} = c_{IZII} = 0.171413$, $c_{IIZI} = c_{IIIZ} = -0.223432$, $c_{ZZII} = 0.168689$, $c_{ZIZI} = c_{IZIZ} = 0.120625$,
$c_{ZIIZ} = c_{IZZI} = 0.165928$, $c_{IIZZ} = 0.174413$, the Hartree--Fock energy is
\par\smallskip\noindent\begin{tabularx}{\linewidth}{@{}*{4}{L}@{}}
\optlabel{A}\,$-1.137284$ Ha & \optlabel{B}\,$-0.097066$ Ha & \optlabel{C}\,$+0.462618$ Ha & \optlabel{D}\,$-1.116759$ Ha
\end{tabularx}
\end{qitem}
\qsep

\needspace{7\baselineskip}
\begin{qitem}{49}
\textbf{Assertion (A):} The notebook's one-parameter UCC ansatz can reach the exact (FCI) ground energy of its H$_2$ Hamiltonian.\\[1mm]
\textbf{Reason (R):} The exact ground state of that Hamiltonian lies in $\mathrm{span}\{\ket{1100}, \ket{0011}\}$, and $\cos\theta\ket{1100} - \sin\theta\ket{0011}$ covers every
real normalised vector of that span (up to sign).
\begin{itemize}[leftmargin=2.6em, labelsep=0.4em, itemsep=2pt, topsep=3pt]
\item[\optlabel{A}] Both Assertion (A) and Reason (R) are true, and R is the correct explanation of A
\item[\optlabel{B}] Both A and R are true, but R is \emph{not} the correct explanation of A
\item[\optlabel{C}] A is true, but R is false
\item[\optlabel{D}] A is false, but R is true
\end{itemize}
\end{qitem}
\qsep

\needspace{7\baselineskip}
\begin{qitem}{50}
For the default Ising chain ($N = 15$, $J = 1$, $h = 0.5$, $H = -J\sum Z_iZ_{i+1} - h\sum X_i$), the energy of $\ket{00\dots0}$ --- close to where the
hardware-efficient ansatz starts, since its initial angles are drawn from $\mathcal N(0, 0.1^2)$ --- is
\par\smallskip\noindent\begin{tabularx}{\linewidth}{@{}*{4}{L}@{}}
\optlabel{A}\,$-15$ & \optlabel{B}\,$-21.5$ & \optlabel{C}\,$-7.5$ & \optlabel{D}\,$-14$
\end{tabularx}
\end{qitem}
\qsep

\needspace{7\baselineskip}
\begin{qitem}{51}
Each of the 45 angles of the default hardware-efficient ansatz enters exactly one $R_Y$ gate. The number of circuit evaluations for one exact
parameter-shift gradient, and for one forward finite-difference gradient (L-BFGS-B), are
\par\smallskip\noindent\begin{tabularx}{\linewidth}{@{}*{4}{L}@{}}
\optlabel{A}\,$90$ and $90$ & \optlabel{B}\,$90$ and $46$ & \optlabel{C}\,$2$ and $46$ & \optlabel{D}\,$45$ and $46$
\end{tabularx}
\end{qitem}
\qsep

\needspace{7\baselineskip}
\begin{qitem}{52}
\textbf{Statement I:} COBYLA needs no gradient evaluations.\\[1mm]
\textbf{Statement II:} The notebook reports the energy of the optimizer's final point \texttt{opt.x} as the VQE energy.
\begin{itemize}[leftmargin=2.6em, labelsep=0.4em, itemsep=2pt, topsep=3pt]
\item[\optlabel{A}] Both Statement I and Statement II are true
\item[\optlabel{B}] Both Statement I and Statement II are false
\item[\optlabel{C}] Statement I is true, but Statement II is false
\item[\optlabel{D}] Statement I is false, but Statement II is true
\end{itemize}
\end{qitem}
\qsep

\needspace{7\baselineskip}
\begin{qitem}{53}
Four VQE runs on the notebook's H$_2$ Hamiltonian (exact expectation values; FCI $= -1.137284$ Ha) report the energies below. Which one is a valid,
chemically accurate result?
\par\smallskip\noindent\begin{tabularx}{\linewidth}{@{}*{4}{L}@{}}
\optlabel{A}\,$-1.11676$ Ha & \optlabel{B}\,$-1.13600$ Ha & \optlabel{C}\,$-1.13500$ Ha & \optlabel{D}\,$-1.13900$ Ha
\end{tabularx}
\end{qitem}
\qsep

\needspace{7\baselineskip}
\begin{qitem}{54}
The notebook's H$_2$ Hamiltonian has the identity, 10 $Z$-type terms and 4 terms $\texttt{XXYY}$, $\texttt{XYYX}$, $\texttt{YXXY}$, $\texttt{YYXX}$. With
\emph{qubit-wise commuting} grouping, the minimum number of measurement settings is ..., and with general commuting grouping ...
\par\smallskip\noindent\begin{tabularx}{\linewidth}{@{}*{4}{L}@{}}
\optlabel{A}\,$5$; $2$ & \optlabel{B}\,$15$; $5$ & \optlabel{C}\,$2$; $1$ & \optlabel{D}\,$4$; $1$
\end{tabularx}
\end{qitem}
\qsep

\needspace{7\baselineskip}
\begin{qitem}{55}
In the previously recorded default run, $E_{\mathrm{VQE}} = -15.0533898284$ and the exact $E_0 = -15.0825179883$. The relative error is
\par\smallskip\noindent\begin{tabularx}{\linewidth}{@{}*{4}{L}@{}}
\optlabel{A}\,$1.93\,\%$ & \optlabel{B}\,$0.193\,\%$ & \optlabel{C}\,$0.0291\,\%$ & \optlabel{D}\,$-0.193\,\%$
\end{tabularx}
\end{qitem}
\qsep

\needspace{7\baselineskip}
\begin{qitem}{56}
The notebook's hardware-efficient ansatz uses only $R_Y$ rotations and CNOTs, so it produces states with \emph{real} amplitudes. Why is this
adequate for its Ising and H$_2$ Hamiltonians?
\begin{itemize}[leftmargin=2.6em, labelsep=0.4em, itemsep=2pt, topsep=3pt]
\item[\optlabel{A}] Complex amplitudes cannot be measured, so they are never needed
\item[\optlabel{B}] Both Hamiltonians are real symmetric matrices, so a ground state can always be chosen real
\item[\optlabel{C}] Real amplitudes always give lower energies than complex ones
\item[\optlabel{D}] CNOT gates generate all complex phases automatically
\end{itemize}
\end{qitem}
\qsep

\needspace{7\baselineskip}
\begin{qitem}{57}
The recorded optimised state of the Ising chain gives $\ket{000\dots0}$ with probability $0.70$ but $\ket{111\dots1}$ (not in the top 8) with nearly zero
probability, although $H$ is symmetric under flipping all spins. The best explanation is
\begin{itemize}[leftmargin=2.6em, labelsep=0.4em, itemsep=2pt, topsep=3pt]
\item[\optlabel{A}] $\ket{111\dots1}$ has a higher energy than $\ket{000\dots0}$
\item[\optlabel{B}] the exact ground state contains only $\ket{000\dots0}$
\item[\optlabel{C}] the simulator has a bias towards $\ket0$
\item[\optlabel{D}] for $h < J$ the two ordered configurations are nearly degenerate; VQE settled on one of them (symmetry breaking), which costs only an energy exponentially small in $N$
\end{itemize}
\end{qitem}
\qsep

\needspace{7\baselineskip}
\begin{qitem}{58}
Match each optimizer offered by the notebook (List-I) with its working principle (List-II).
\par\smallskip\noindent{\renewcommand{\arraystretch}{1.35}%
\begin{tabularx}{\linewidth}{@{}c>{\raggedright\arraybackslash}p{0.33\linewidth}@{\hspace{1.2em}}cL@{}}
\toprule
\multicolumn{2}{@{}l}{\textbf{List-I}} & \multicolumn{2}{l@{}}{\textbf{List-II}} \\ \midrule
\textbf{(P)} & COBYLA & \textbf{(1)} & successive line searches along a set of conjugate directions, derivative-free \\
\textbf{(Q)} & Nelder--Mead & \textbf{(2)} & quasi-Newton method with a limited-memory Hessian estimate, using gradients (finite differences here) and box bounds \\
\textbf{(R)} & Powell & \textbf{(3)} & linear models of the objective inside a shrinking trust region, derivative-free \\
\textbf{(S)} & L-BFGS-B & \textbf{(4)} & a simplex of $n+1$ points updated by reflection, expansion and contraction \\
\bottomrule
\end{tabularx}}
\par\smallskip\noindent\begin{tabularx}{\linewidth}{@{}*{2}{L}@{}}
\optlabel{A}\,P$\to$3, Q$\to$4, R$\to$1, S$\to$2 & \optlabel{B}\,P$\to$3, Q$\to$1, R$\to$4, S$\to$2 \\[3pt]
\optlabel{C}\,P$\to$4, Q$\to$3, R$\to$1, S$\to$2 & \optlabel{D}\,P$\to$3, Q$\to$4, R$\to$2, S$\to$1
\end{tabularx}
\end{qitem}
\qsep

\needspace{20\baselineskip}
\begin{passagebox}{Paragraph for Questions 59 and 60}
In the notebook's H$_2$ example the UCC state $\ket{\psi(\theta)} = \cos\theta\ket{1100} - \sin\theta\ket{0011}$ only involves two basis states. Within
this two-dimensional space the Hamiltonian is the real $2\times2$ matrix
\[ \begin{pmatrix} H_{11} & H_{12}\\ H_{12} & H_{22}\end{pmatrix} = \begin{pmatrix} -1.116759 & 0.181210\\ 0.181210 & 0.462618\end{pmatrix}\ \text{Ha}, \]
with $H_{11} = \braket{1100|H|1100}$, $H_{22} = \braket{0011|H|0011}$, $H_{12} = \braket{1100|H|0011}$ (computed from the notebook's 15 Pauli terms), so that
$E(\theta) = H_{11}\cos^2\theta + H_{22}\sin^2\theta - H_{12}\sin2\theta$.
\end{passagebox}
\needspace{7\baselineskip}
\begin{qitem}{59}
The lowest eigenvalue of the $2\times2$ matrix --- the best energy the UCC ansatz can reach --- is
\par\smallskip\noindent\begin{tabularx}{\linewidth}{@{}*{4}{L}@{}}
\optlabel{A}\,$-1.116759$ Ha & \optlabel{B}\,$+0.483143$ Ha & \optlabel{C}\,$-0.654141$ Ha & \optlabel{D}\,$-1.137284$ Ha
\end{tabularx}
\end{qitem}
\qsep

\needspace{7\baselineskip}
\begin{qitem}{60}
The optimal angle $\theta^*$ and the probability of measuring $\ket{0011}$ in the optimal state are
\par\smallskip\noindent\begin{tabularx}{\linewidth}{@{}*{2}{L}@{}}
\optlabel{A}\,$\theta^* \approx 0.2295$ rad and $0.0518$ & \optlabel{B}\,$\theta^* \approx 0.1128$ rad and $0.0127$ \\[3pt]
\optlabel{C}\,$\theta^* \approx -0.1128$ rad and $0.0127$ & \optlabel{D}\,$\theta^* \approx 0.1128$ rad and $0.1125$
\end{tabularx}
\end{qitem}
\qsep

\newpage

\begin{sectionbanner}SECTION D \enspace$\cdot$\enspace Advanced VQE and Implementation \hfill Q.61--80\end{sectionbanner}

{\small \textbf{Syllabus:} Jordan--Wigner and Bravyi--Kitaev mappings, UCC, molecular scaling, shots, noise and error mitigation, backends, exact references, benchmarking and interpretation of recorded results.\\ \textbf{Formats:} Single-correct MCQ (Q.61--78); Matching lists (Q.64--68); Assertion--Reason (Q.71); Statement I / II (Q.74); Paragraph based (Q.79--80). Each question: $+3$ / $0$ / $-1$.}
\par\medskip

\needspace{7\baselineskip}
\begin{qitem}{61}
With the Jordan--Wigner transformation $a_j^\dagger = \big(\prod_{k<j}Z_k\big)\,\tfrac12(X_j - iY_j)$ (qubit $\ket1$ = occupied), the creation operator $a_2^\dagger$ on
4 spin-orbitals becomes
\par\smallskip\noindent\begin{tabularx}{\linewidth}{@{}*{2}{L}@{}}
\optlabel{A}\,$\tfrac12\,Z_0Z_1(X_2 + iY_2)$ & \optlabel{B}\,$\tfrac12(X_2 - iY_2)$ \\[3pt]
\optlabel{C}\,$\tfrac12\,Z_3(X_2 - iY_2)$ & \optlabel{D}\,$\tfrac12\,Z_0Z_1(X_2 - iY_2)$
\end{tabularx}
\end{qitem}
\qsep

\needspace{7\baselineskip}
\begin{qitem}{62}
Under Jordan--Wigner, the hopping term $a_0^\dagger a_1 + a_1^\dagger a_0$ between neighbouring spin-orbitals 0 and 1 becomes
\par\smallskip\noindent\begin{tabularx}{\linewidth}{@{}*{2}{L}@{}}
\optlabel{A}\,$X_0X_1$ & \optlabel{B}\,$\tfrac12(X_0X_1 + Y_0Y_1)$ \\[3pt]
\optlabel{C}\,$\tfrac12(Z_0 + Z_1)$ & \optlabel{D}\,$\tfrac12(X_0X_1 - Y_0Y_1)$
\end{tabularx}
\end{qitem}
\qsep

\needspace{7\baselineskip}
\begin{qitem}{63}
The Pauli weight (number of non-identity factors) of the image of $a_7^\dagger$ on 8 spin-orbitals is ... under Jordan--Wigner and ... under Bravyi--Kitaev.
\par\smallskip\noindent\begin{tabularx}{\linewidth}{@{}*{4}{L}@{}}
\optlabel{A}\,$8$; $4$ & \optlabel{B}\,$8$; $8$ & \optlabel{C}\,$1$; $4$ & \optlabel{D}\,$4$; $8$
\end{tabularx}
\end{qitem}
\qsep

\needspace{7\baselineskip}
\begin{qitem}{64}
Match each fermion-to-qubit encoding (List-I) with its characteristic property (List-II).
\par\smallskip\noindent{\renewcommand{\arraystretch}{1.35}%
\begin{tabularx}{\linewidth}{@{}c>{\raggedright\arraybackslash}p{0.33\linewidth}@{\hspace{1.2em}}cL@{}}
\toprule
\multicolumn{2}{@{}l}{\textbf{List-I}} & \multicolumn{2}{l@{}}{\textbf{List-II}} \\ \midrule
\textbf{(P)} & Jordan--Wigner & \textbf{(1)} & qubit $j$ stores the parity $n_0 \oplus\dots\oplus n_j$; total parity and $S_z$-parity sit on single qubits, so two qubits can be removed \\
\textbf{(Q)} & Bravyi--Kitaev & \textbf{(2)} & symmetries ($\mathbb Z_2$ stabilisers commuting with $H$) are used to remove qubits from any encoding \\
\textbf{(R)} & parity mapping & \textbf{(3)} & qubit $j$ stores the occupation $n_j$; strings of length $O(N)$ \\
\textbf{(S)} & qubit tapering & \textbf{(4)} & qubits store partial sums of occupations on a binary tree; strings of length $O(\log N)$ \\
\bottomrule
\end{tabularx}}
\par\smallskip\noindent\begin{tabularx}{\linewidth}{@{}*{2}{L}@{}}
\optlabel{A}\,P$\to$4, Q$\to$3, R$\to$1, S$\to$2 & \optlabel{B}\,P$\to$3, Q$\to$4, R$\to$1, S$\to$2 \\[3pt]
\optlabel{C}\,P$\to$3, Q$\to$1, R$\to$4, S$\to$2 & \optlabel{D}\,P$\to$3, Q$\to$4, R$\to$2, S$\to$1
\end{tabularx}
\end{qitem}
\qsep

\needspace{7\baselineskip}
\begin{qitem}{65}
For 4 modes, the Bravyi--Kitaev qubit values are $b = Bn \bmod 2$ with $B = \begin{psmallmatrix}1&0&0&0\\1&1&0&0\\0&0&1&0\\1&1&1&1\end{psmallmatrix}$ and $n$ the occupation vector.
The Hartree--Fock occupation $n = (1, 1, 0, 0)$ is encoded as
\par\smallskip\noindent\begin{tabularx}{\linewidth}{@{}*{4}{L}@{}}
\optlabel{A}\,$\ket{0100}$ & \optlabel{B}\,$\ket{1100}$ & \optlabel{C}\,$\ket{1101}$ & \optlabel{D}\,$\ket{1000}$
\end{tabularx}
\end{qitem}
\qsep

\needspace{7\baselineskip}
\begin{qitem}{66}
Counting excitations in spin-orbitals without spin restrictions, LiH in STO-3G (12 spin-orbitals, 4 electrons) has how many single and double
excitations from the Hartree--Fock determinant?
\par\smallskip\noindent\begin{tabularx}{\linewidth}{@{}*{2}{L}@{}}
\optlabel{A}\,$32$ singles and $168$ doubles & \optlabel{B}\,$32$ singles and $1024$ doubles \\[3pt]
\optlabel{C}\,$8$ singles and $28$ doubles & \optlabel{D}\,$48$ singles and $168$ doubles
\end{tabularx}
\end{qitem}
\qsep

\needspace{7\baselineskip}
\begin{qitem}{67}
Under Jordan--Wigner, the anti-Hermitian double-excitation generator $a_3^\dagger a_2^\dagger a_1a_0 - a_0^\dagger a_1^\dagger a_2a_3$ becomes a sum of how many Pauli strings,
and what does the notebook's UCC kernel use?
\begin{itemize}[leftmargin=2.6em, labelsep=0.4em, itemsep=2pt, topsep=3pt]
\item[\optlabel{A}] 4 strings with coefficients $\pm\tfrac12$; the notebook keeps $\texttt{XXYY}$
\item[\optlabel{B}] 1 string; the notebook uses all of it
\item[\optlabel{C}] 8 strings with coefficients $\pm i/8$; the notebook keeps the single string $\texttt{XXXY}$
\item[\optlabel{D}] 16 strings; the notebook keeps 8
\end{itemize}
\end{qitem}
\qsep

\needspace{7\baselineskip}
\begin{qitem}{68}
Match each practical VQE difficulty (List-I) with a standard remedy (List-II).
\par\smallskip\noindent{\renewcommand{\arraystretch}{1.35}%
\begin{tabularx}{\linewidth}{@{}c>{\raggedright\arraybackslash}p{0.33\linewidth}@{\hspace{1.2em}}cL@{}}
\toprule
\multicolumn{2}{@{}l}{\textbf{List-I}} & \multicolumn{2}{l@{}}{\textbf{List-II}} \\ \midrule
\textbf{(P)} & barren plateaus (vanishing gradients for deep random circuits) & \textbf{(1)} & error mitigation, e.g.\ zero-noise extrapolation \\
\textbf{(Q)} & shot noise in energy estimates & \textbf{(2)} & chemically motivated ansatz, local cost functions, careful (e.g.\ Hartree--Fock) initialisation \\
\textbf{(R)} & gate errors on hardware & \textbf{(3)} & several restarts / good initial parameters \\
\textbf{(S)} & local minima of $E(\boldsymbol\theta)$ & \textbf{(4)} & measurement grouping and more shots for terms with large $|c_k|$ \\
\bottomrule
\end{tabularx}}
\par\smallskip\noindent\begin{tabularx}{\linewidth}{@{}*{2}{L}@{}}
\optlabel{A}\,P$\to$4, Q$\to$2, R$\to$1, S$\to$3 & \optlabel{B}\,P$\to$2, Q$\to$4, R$\to$1, S$\to$3 \\[3pt]
\optlabel{C}\,P$\to$2, Q$\to$1, R$\to$4, S$\to$3 & \optlabel{D}\,P$\to$2, Q$\to$4, R$\to$3, S$\to$1
\end{tabularx}
\end{qitem}
\qsep

\needspace{7\baselineskip}
\begin{qitem}{69}
An energy estimator has a per-shot standard deviation of $0.1$ Ha. How many shots are needed for a standard error of $1.6$ mHa (chemical accuracy)?
\par\smallskip\noindent\begin{tabularx}{\linewidth}{@{}*{4}{L}@{}}
\optlabel{A}\,about $1.6\times10^{5}$ & \optlabel{B}\,about $63$ & \optlabel{C}\,about $3.9\times10^3$ & \optlabel{D}\,about $2.6\times10^{-4}$
\end{tabularx}
\end{qitem}
\qsep

\needspace{7\baselineskip}
\begin{qitem}{70}
Under complete depolarisation the state becomes $\rho = I/2^n$. For the notebook's 4-qubit H$_2$ Hamiltonian the measured energy would then be
\par\smallskip\noindent\begin{tabularx}{\linewidth}{@{}*{2}{L}@{}}
\optlabel{A}\,$-0.0971$ Ha, the coefficient of $\texttt{IIII}$ & \optlabel{B}\,$-1.1168$ Ha \\[3pt]
\optlabel{C}\,$-1.1373$ Ha & \optlabel{D}\,$0$ Ha
\end{tabularx}
\end{qitem}
\qsep

\needspace{7\baselineskip}
\begin{qitem}{71}
\textbf{Assertion (A):} With a finite number of shots, an estimated VQE energy can never fall below $E_0$.\\[1mm]
\textbf{Reason (R):} For every density matrix $\rho$ (pure or mixed, noisy or not), $\operatorname{tr}(\rho H) \ge E_0$.
\begin{itemize}[leftmargin=2.6em, labelsep=0.4em, itemsep=2pt, topsep=3pt]
\item[\optlabel{A}] Both Assertion (A) and Reason (R) are true, and R is the correct explanation of A
\item[\optlabel{B}] Both A and R are true, but R is \emph{not} the correct explanation of A
\item[\optlabel{C}] A is true, but R is false
\item[\optlabel{D}] A is false, but R is true
\end{itemize}
\end{qitem}
\qsep

\needspace{7\baselineskip}
\begin{qitem}{72}
Zero-noise extrapolation: an energy measured at noise scale factors $\lambda = 1$ and $\lambda = 3$ gives $-1.10$ Ha and $-1.00$ Ha. A linear extrapolation to
$\lambda = 0$ gives
\par\smallskip\noindent\begin{tabularx}{\linewidth}{@{}*{4}{L}@{}}
\optlabel{A}\,$-1.10$ Ha & \optlabel{B}\,$-1.15$ Ha & \optlabel{C}\,$-1.05$ Ha & \optlabel{D}\,$-1.20$ Ha
\end{tabularx}
\end{qitem}
\qsep

\needspace{7\baselineskip}
\begin{qitem}{73}
\texttt{VQE.ipynb} is run with \texttt{BACKEND = "gpu"} on a machine where \texttt{cudaq.num\_available\_gpus()} returns 0. What happens?
\begin{itemize}[leftmargin=2.6em, labelsep=0.4em, itemsep=2pt, topsep=3pt]
\item[\optlabel{A}] The notebook runs on the GPU emulator \texttt{tensornet} instead
\item[\optlabel{B}] CUDA-Q silently falls back to \texttt{qpp-cpu}
\item[\optlabel{C}] \texttt{resolve\_backend} raises a \texttt{RuntimeError} explaining that no NVIDIA GPU is visible; no GPU target is selected
\item[\optlabel{D}] The run starts and crashes later with an out-of-memory error
\end{itemize}
\end{qitem}
\qsep

\needspace{7\baselineskip}
\begin{qitem}{74}
\textbf{Statement I:} For the Ising chain, the notebook computes its reference energy by sparse diagonalisation for every $N$.\\[1mm]
\textbf{Statement II:} For \texttt{"H2"} and \texttt{"custom"} Hamiltonians, an exact reference is computed only up to \texttt{REFERENCE\_MAX\_QUBITS = 18} qubits.
\begin{itemize}[leftmargin=2.6em, labelsep=0.4em, itemsep=2pt, topsep=3pt]
\item[\optlabel{A}] Both Statement I and Statement II are true
\item[\optlabel{B}] Both Statement I and Statement II are false
\item[\optlabel{C}] Statement I is true, but Statement II is false
\item[\optlabel{D}] Statement I is false, but Statement II is true
\end{itemize}
\end{qitem}
\qsep

\needspace{7\baselineskip}
\begin{qitem}{75}
The notebook's exact Ising reference maps the chain to free fermions and computes $E_0 = -\tfrac12\sum|\lambda(iA)|$ for a real antisymmetric matrix $A$. For
$N = 100$ qubits, $A$ has size
\par\smallskip\noindent\begin{tabularx}{\linewidth}{@{}*{4}{L}@{}}
\optlabel{A}\,$4\times4$ & \optlabel{B}\,$200\times200$ & \optlabel{C}\,$100\times100$ & \optlabel{D}\,$2^{100}\times2^{100}$
\end{tabularx}
\end{qitem}
\qsep

\needspace{7\baselineskip}
\begin{qitem}{76}
In the recorded run the quantum execution time was $69.10$ s for $300$ energy evaluations, while the notebook's pre-run estimate was $\sim102.9$ s.
The average time per evaluation, and the ratio estimate/measured cell time ($70.1$ s), are about
\par\smallskip\noindent\begin{tabularx}{\linewidth}{@{}*{2}{L}@{}}
\optlabel{A}\,$0.0032$ s and $1.5$ & \optlabel{B}\,$69$ s and $1.5$ \\[3pt]
\optlabel{C}\,$0.23$ s and $0.68$ & \optlabel{D}\,$0.23$ s and $1.5$
\end{tabularx}
\end{qitem}
\qsep

\needspace{7\baselineskip}
\begin{qitem}{77}
For LiH (STO-3G, 12 spin-orbitals under Jordan--Wigner), freezing the doubly occupied Li $1s$ core orbital removes how many qubits?
\par\smallskip\noindent\begin{tabularx}{\linewidth}{@{}*{2}{L}@{}}
\optlabel{A}\,$2$, leaving $10$ qubits & \optlabel{B}\,$0$; frozen orbitals still need qubits \\[3pt]
\optlabel{C}\,$4$, leaving $8$ qubits & \optlabel{D}\,$1$, leaving $11$ qubits
\end{tabularx}
\end{qitem}
\qsep

\needspace{7\baselineskip}
\begin{qitem}{78}
Generated with PySCF and OpenFermion (STO-3G, Jordan--Wigner, bond length 1.595 \AA), the LiH qubit Hamiltonian has how many Pauli terms, compared with
15 for H$_2$?
\par\smallskip\noindent\begin{tabularx}{\linewidth}{@{}*{4}{L}@{}}
\optlabel{A}\,$4096$ & \optlabel{B}\,$144$ & \optlabel{C}\,$631$ & \optlabel{D}\,$30$
\end{tabularx}
\end{qitem}
\qsep

\needspace{20\baselineskip}
\begin{passagebox}{Paragraph for Questions 79 and 80}
\emph{Previously recorded default run of \texttt{VQE.ipynb}} (CUDA-Q 0.16.0, target \texttt{qpp-cpu}, 4 vCPU): Ising chain with $N = 15$, $J = 1$,
$h = 0.5$ (29 Pauli terms); hardware-efficient ansatz with 2 layers (45 parameters, 73 gates); COBYLA stopped with ``the objective function has
been evaluated MAXFUN times'' after 300 evaluations; $E_{\mathrm{VQE}} = -15.0533898284$, exact $E_0 = -15.0825179883$; quantum execution time
$69.10$ s; the optimised state gave $\ket{000\dots0}$ with probability $0.7005$. The product state $\ket{000\dots0}$ itself has energy $-14$.
\end{passagebox}
\needspace{7\baselineskip}
\begin{qitem}{79}
Measured from the product state $\ket{000\dots0}$ ($E = -14$), what fraction of the remaining energy gap to $E_0$ did VQE recover?
\par\smallskip\noindent\begin{tabularx}{\linewidth}{@{}*{4}{L}@{}}
\optlabel{A}\,$97.3\,\%$ & \optlabel{B}\,$99.8\,\%$ & \optlabel{C}\,$70.1\,\%$ & \optlabel{D}\,$2.7\,\%$
\end{tabularx}
\end{qitem}
\qsep

\needspace{7\baselineskip}
\begin{qitem}{80}
Which change is most likely to reduce the remaining error of this run?
\begin{itemize}[leftmargin=2.6em, labelsep=0.4em, itemsep=2pt, topsep=3pt]
\item[\optlabel{A}] More final sampling shots
\item[\optlabel{B}] Switching from the free-fermion reference to sparse diagonalisation
\item[\optlabel{C}] More evaluations (the budget was exhausted) and, if needed, more ansatz layers
\item[\optlabel{D}] Reducing \texttt{MAX\_EVALUATIONS} to stop earlier
\end{itemize}
\end{qitem}
\qsep

\newpage

\end{document}